%------------------------------------------------------------------------------%
\documentclass[a4paper,12pt,fleqn]{article}

\usepackage{calculo_varias_variaveis-1}

\ShowAnswers

%------------------------------------------------------------------------------%
\begin{document}

\makeHeader{CFVVI}{Prova 5 -- Turma 2}{17/12/2025}

% 01
%------------------------------------------------------------------------------%
\exercise{25\%}
Considere a equação em coordenadas polares
\(
    4r = \left(1+\cos\theta\right)^{-1},
\)
transforme-a em coordenadas cartesianas,
identifique a cônica obtida e esboce seu gráfico.

\begin{questiononly}
  \bigskip
  \hfill\includegraphics{axis_xy}
\end{questiononly}
\clearpagequestiononly

\begin{answer}
\begin{multicols}{2}
\begin{align*}
    4r                       & = \left(1+\cos\theta\right)^{-1} \\
    4r                       & = \frac{1}{1+\cos\theta}         \\
    4r(1+\cos\theta)         & = 1                              \\
    4r+4r\cos\theta          & = 1                              \\
    4r+4x                    & = 1                              \\
    4r                       & = 1 - 4x                         \\
    16r^2                    & = (1 - 4x)^2                     \\
    16\left(x^2 + y^2\right) & = 1 -8x + 16x^2                  \\
    16y^2                    & = 1 -8x                          \\
    8x                       & = 1 - 16y^2                      \\
    x                        & = \frac{1}{8} -2y^2
\end{align*}

\columnbreak
A cônica é uma parábola
\bigskip

\includegraphics{T2_A_conica}
\end{multicols}
\end{answer}

% 02
%------------------------------------------------------------------------------%
\exercise{25\%}
Determine todas as raízes cúbicas de
\(
    z = 4\sqrt{3}-4i
\)

\clearpagequestiononly

\begin{answer}
Módulo de $z$
\[
    \abs{z}
    = \sqrt{\left(4\sqrt{3}\right)^2+(-4)^2}
    = \sqrt{48+16}
    = 8
\]
Para encontrar o argumento \(\varphi = \arg(z)\),
usamos o seno e o cosseno do âlgulo
\[
    \sin(\varphi)
    = \frac{\Im(z)}{\abs{z}}
    = \frac{-4}{8}
    = \frac{-1}{2}
\]
\[
    \cos(\varphi)
    = \frac{\Re(z)}{\abs{z}}
    = \frac{4\sqrt{3}}{8}
    = \frac{\sqrt{3}}{2}
\]
portanto \(\varphi=\frac{12n -1}{6}\pi\), escolhemos \(\varphi=\frac{-\pi}{6}\)
\[
    z = 8
    \left(
        \cos\left(\frac{-\pi}{6}\right) +
        i\sin\left(\frac{-\pi}{6}\right)
    \right)
\]

Pela segunda fórmula de Moivre, as raízes cúbicas são
\[
    u_k=\sqrt[3]{8}
    \left(
          \cos\left(\frac{-\sfrac{\pi}{6}+2k\pi}{3}\right)
        +i\sin\left(\frac{-\sfrac{\pi}{6}+2k\pi}{3}\right)
    \right)
    \qquad k = 0, 1, 2
\]

Como $\sqrt[3]{8}=2$, obtemos
\[
    u_0 = 2 \left(
                \cos\left(-\frac{\pi}{18}\right)
                +i\sin\left(-\frac{\pi}{18}\right)
            \right)
\]
\[
    u_1 = 2 \left(
                \cos\left(\frac{11\pi}{18}\right)
                +i\sin\left(\frac{11\pi}{18}\right)
            \right)
\]
\[
    u_2 = 2 \left(
                \cos\left(\frac{23\pi}{18}\right)
                +i\sin\left(\frac{23\pi}{18}\right)
            \right)
\]
\end{answer}


% 03
%------------------------------------------------------------------------------%
\exercise{50\%}
Determine os pontos candidatos a extremos de
\(
    x^2 + \ln y
\)
com $y>0$, sujeita à restrição
\(
    2x + y = 3
\)

\begin{answer}
\begin{multicols}{2}
Queremos os extremos de
\[
    f(x, y) = x^2+\ln y \qquad y > 0
\]
sujeita à restrição
\[
    g(x, y) = 2x + y = 3
\]
Gradientes
\[
    \nabla f = \vetor{2x}{\sfrac{1}{y}}
    \qquad
    \nabla g = \vetor{2}{1}
\]

Pelo método dos multiplicadores de Lagrange
\begin{align*}
    2x          & = 2\lambda \\
    \frac{1}{y} & = \lambda \\
    2x + y      & = 3
\end{align*}
Da primeira equação temos que \(x=\lambda\), \\
substituindo na segunda temos
\[
    \frac{1}{y} = \lambda = x
    \qquad
    y = \frac{1}{x}
\]
Substituindo na restrição
\begin{align*}
    2x + y           & = 3  \\
    2x + \frac{1}{x} & = 3  \\
    2x^2 + 1         & = 3x \\
    2x^2 - 3x + 1    & = 0
\end{align*}
Resolvendo
\[
    \Delta
    = b^2 - 4ac
    = (-3)^2 - 4 \times 2 \times 1
    = 9 - 8
    = 1
\]
\[
    x
    = \frac{-b\pm\sqrt{\Delta}}{2a}
    = \frac{3\pm 1}{4}
\]

\smallskip
Para \(x=\frac{3+1}{4}=1\) temos \(y=1\)

\smallskip
Para \(x=\frac{3-1}{4}=\frac{1}{2}\) temos \(y=2\)

\smallskip
Os pontos candidatos a extremos são
\[
  P_1 = \left(1, 1\right)
  \quad \text{ e } \quad
  P_2 = \left(\frac{1}{2}, 2\right)
\]
\end{multicols}
\end{answer}

%------------------------------------------------------------------------------%
\end{document}
%------------------------------------------------------------------------------%
